\section{Draft Results}

This draft section summarizes the current state of the prototype. The main takeaway is that the deterministic RF heuristic already exposes a meaningful borderline regime, while the current bounded LLM adjudicator is instrumented and testable but does not yet improve performance on that regime. This is still a useful result because it gives us a concrete, interpretable baseline and shows exactly where the LLM needs improvement.

\subsection{Evaluation Setup}

We generated a weak-to-strong synthetic attack sweep over 82 scenarios spanning four attack families: sensor foil, gateway bias, random noise, and packet drop. Attack strengths were varied so that the evaluation included both easy cases and borderline cases near the heuristic decision boundary. The current agent was evaluated in two modes:

\begin{itemize}
    \item \textbf{Heuristic-only}: deterministic RF scoring without any LLM intervention.
    \item \textbf{Bounded adjudication}: the same heuristic detector, with the LLM invoked only on scenarios flagged as ambiguous.
\end{itemize}

We report overall accuracy, ambiguous-case accuracy, false positive rate on scenarios labeled as benign, and LLM invocation rate. The purpose of this analysis is not to show that the LLM universally outperforms heuristics, but to identify whether there is an interpretable ambiguity regime where escalation is justified.

\subsection{Headline Metrics}

\begin{figure}[t]
    \centering
    \includegraphics[width=\columnwidth]{../outputs/advisor_demo/mode_metrics.png}
    \caption{Summary comparison between heuristic-only inference and bounded LLM adjudication. The current prototype shows that the LLM is only invoked on a small fraction of scenarios, but it does not yet improve accuracy.}
    \label{fig:mode-metrics}
\end{figure}

The current headline metrics are:

\begin{itemize}
    \item Heuristic-only overall accuracy: 0.695
    \item Heuristic-only ambiguous-case accuracy: 0.286
    \item Adjudication-mode overall accuracy: 0.683
    \item Adjudication-mode ambiguous-case accuracy: 0.238
    \item Adjudication-mode LLM invocation rate: 0.098
\end{itemize}

The most important observation is that the evaluation is no longer trivial. On strong attacks the heuristic detector performs well, but once attack strength drops into the weak-to-medium regime, performance falls substantially. This is the region where a higher-level adjudication mechanism could matter. At the moment, however, the current LLM implementation does not yet outperform the deterministic baseline.

\subsection{Where Ambiguity Appears}

\begin{figure*}[t]
    \centering
    \includegraphics[width=0.92\textwidth]{../outputs/advisor_demo/ambiguity_vs_severity.png}
    \caption{Ambiguity and accuracy as a function of attack severity. The plots show that the detector is confident on strong attacks, while ambiguity is concentrated in the weak-to-medium regime.}
    \label{fig:ambiguity-severity}
\end{figure*}

Figure~\ref{fig:ambiguity-severity} shows where ambiguity arises across attack families. The heuristic raises ambiguity flags primarily when attack severity is strong enough to disturb the clean baseline, but not strong enough to match one attack signature decisively. This is the exact operating region where a supervisor-style model is most plausible.

In the current sweep, 21 out of 82 scenarios were flagged as borderline by the heuristic. The most common ambiguity triggers were:

\begin{itemize}
    \item low heuristic confidence
    \item no strong rule firing
    \item anomaly score above the clean threshold without a decisive attack signature
    \item borderline noise signatures in the variance and trilateration residual signals
\end{itemize}

This is useful from a systems perspective because it suggests the escalation policy is at least targeting a non-trivial regime rather than calling the LLM on every case.

\subsection{Interpretable Case Studies}

\begin{figure*}[t]
    \centering
    \includegraphics[width=0.92\textwidth]{../outputs/advisor_demo/case_study_heatmaps.png}
    \caption{Per-sensor / per-gateway RSSI shift patterns for three borderline cases. These examples make it visually clear why some cases are harder than strong attacks.}
    \label{fig:case-heatmaps}
\end{figure*}

\begin{figure*}[t]
    \centering
    \includegraphics[width=0.82\textwidth]{../outputs/advisor_demo/localization_case_study.png}
    \caption{Geometry-based cross-checks for representative cases. Localization error and residual behavior provide a physics-grounded complement to RSSI-only inspection.}
    \label{fig:localization-cases}
\end{figure*}

The case studies help explain what the heuristic and LLM are currently reasoning over:

\begin{itemize}
    \item \textbf{Borderline sensor foil}: the anomaly is coherent across all gateways for one sensor, but the shift magnitude is weak enough that the heuristic is not fully decisive.
    \item \textbf{Borderline packet drop}: only one sensor-gateway relation shows a weak packet-retention change, which can look like mild loss rather than a clear attack.
    \item \textbf{Borderline random noise}: the disturbance appears more as variance inflation than as a strong mean RSSI shift, making it harder to separate from benign RF instability.
\end{itemize}

Figure~\ref{fig:case-heatmaps} shows the corresponding RSSI delta patterns, while Figure~\ref{fig:localization-cases} shows that the trilateration-based geometry checks provide an additional, more physical signal that is sometimes complementary to the raw RSSI evidence.

\subsection{Current Interpretation}

The current result is best framed as follows:

\begin{itemize}
    \item The heuristic detector is a meaningful and interpretable baseline.
    \item There is a real ambiguous regime where RF evidence is not cleanly separable.
    \item The current bounded LLM adjudicator is operational, selective, and instrumented for failure analysis.
    \item The present LLM prompting and faithfulness constraints are not yet strong enough to improve accuracy.
\end{itemize}

This is still a useful milestone because it gives us a concrete baseline story: the project already shows where RF-only heuristics fail, and it now has a measurable interface for testing whether an LLM can help in those failures without replacing the underlying measurement system.
